\documentclass[10pt,a4paper]{amsart}
\usepackage[margin=19mm]{geometry}
\usepackage{amsmath,amssymb,mathtools}
\usepackage[scr=rsfs,cal=euler]{mathalfa}
\usepackage{xfrac}
\usepackage[unicode,hidelinks,bookmarks=false]{hyperref}
\newcommand{\ie}{i.e.\ }
\newcommand{\cf}{cf.\ }
\newcommand{\resp}{respectively\ }
\newcommand{\Subsection}{\subsection}
\providecommand{\nobreakdash}{\nobreak\hbox{-}}
\setlength{\emergencystretch}{2em}
\multlinegap0pt
\usepackage{amssymb}
\usepackage{xcolor}
\usepackage{enumerate}
\usepackage{siunitx}
\usepackage{algorithm}\usepackage{algpseudocode}
\newtheorem{cor}{Corollary}
\usepackage{cleveref}
\crefname{cor}{Corollary}{Corollarys}
\title{The search for exotic knot traces}
\author{Marc Kegel, Jonathan Spreer}
\date{Source: arXiv:2603.22438v2 (CC BY 4.0)}
\begin{document}
\maketitle

\noindent\emph{Source and licence.} The search for exotic knot traces. Version arXiv:2603.22438v2 (CC BY 4.0), distributed by arXiv under the Creative Commons Attribution 4.0 International licence, http://creativecommons.org/licenses/by/4.0/, as recorded in the arXiv licence metadata for this exact version and shown on the abstract page https://arxiv.org/abs/2603.22438v2. Full licence notice: \texttt{licenses/arxiv-2603.22438-CC-BY-4.0.txt}.\par\smallskip

\noindent\textit{Editorial scope.} Counted unit: the article's cor cor:ribbon (source0.tex lines 641--643). The article's complete original proof of it is delivered with it (source0.tex lines 645--653, one \texttt{proof} environment of 9 lines). No mathematics is written, repaired or re-derived: the statement and the proof are the article's own lines, in the article's own notation, and no retained line is substituted or re-flowed.  Retained uncounted as the source-local context the proof needs: lines 634--639.  One declared adaptation: exactly 1 ancillary strings inside the retained spans are omitted (image-inclusion commands and, where the source repeats a label, the repeated label definitions) and no other line differs from the source; the figure environments, with their placement, captions and labels, are kept, and the package records the omitted strings in fidelity receipts bound to the affected bodies.  No retained line contains a citation, so the wrapper carries no bibliography and no bibliography entry.  The retained lines contain the article's own diagram code, which the wrapper typesets with the standard tikz packages; no image file is delivered and the package declares no asset dependency.  Editorial formatting only, in addition: an AMS \texttt{amsart} preamble that restates the article's own declarations and macros used by the retained lines, namely the environments \texttt{cor} on separate counters, together with the standard packages those lines need.  The retained environments print in this document as Corollary 1 in order of appearance, whereas the article prints the same items with the numbering of its own full document.  The complete native source of the article as distributed in the arXiv e-print for this exact version is bundled unchanged as \texttt{sources/197019/source0.tex}.
 \begin{figure}[htbp]
     \centering
       
     \caption{Attaching a ribbon band to create a concordant knot.}
     \label{fig:ribbon_band}
 \end{figure}

 \begin{cor}\label{cor:ribbon}
     If $K'$ is obtained from $K$ by a ribbon band attachment, then $K$ and $K'$ are concordant. In particular, they have the same $4$-genus.
 \end{cor}

 \begin{proof}
    On the right of \Cref{fig:ribbon_band}, we see the ribbon band $R$ in a $3$-ball $B$. In the figure, we also observe that there are two embedded disks $D_1$ and $D_2$, such that the boundary of each disk consists of two arcs, one arc on the boundary of $B$ and the other coinciding with a component of $R$. These two disks intersect only in their interiors in two ribbon singularities.

    In the product $S^3 \times [0,1]$, we can resolve these two ribbon singularities and obtain two embedded disks. We attach these to the trivial product cobordism 
\[
K \times [0,1] \subset S^3 \times [0,1]
\]
and obtain a ribbon concordance between $K$ and $K'$.
 \end{proof}

\end{document}
